\chapter{Extensions survivantes, monodromie et dynamique du quotient}
\label{chap:extensiones-dinamica-cociente}
\index{survie!extension globale}
\index{monodromie!retour nonadique}
\index{mesure!catalogue d’émissions}

Le langage des chemins nécessite une précision décisive. Un chemin
d'APP, une émission visible, une histoire enrichie de TPK et une section
globale ne sont pas quatre noms pour le même objet. Ils constituent quatre niveaux
consécutifs d'une même construction. Le chemin conserve l'ordre
positionnel ; l'émission conserve une lecture finie ; l'histoire ajoute la signature,
la retenue, l'orientation, la frontière et la mémoire ; la section globale exige que tous
ses préfixes survivent simultanément à toute profondeur.

Cette hiérarchie permet de séparer deux questions qui étaient restées
entremêlées. La première est logique : sous quelles hypothèses un préfixe fini
appartient réellement à une extension infinie compatible. La seconde est
dynamique : ce que signifie transporter une seule histoire et ce que signifie
transférer simultanément toutes ses continuations compatibles. La
distinction est essentielle dans APP\@. Une trajectoire déterminée mène un état à
un autre ; un transfert de chemins mène une distribution d'extensions à
une autre distribution.

Cette section résout le problème fini de profondeur six. La
chronologie déterministe publiée possède un retour visible et ne peut pas sélectionner
une probabilité unique. En revanche, le transport biaisé par la phase nonadique,
lorsque chaque canal actif est interprété comme le renouvellement normalisé de toute
sa fibre de continuations compatibles, possède un retour de neuf pas
uniforme sur les (468) émissions. Les poids (1/3) cessent alors d'être
un choix extérieur : ils proviennent des trois apparitions de chaque signature
(+1,-1,0) dans le cycle nonadique orienté. La normalisation au sein de chaque fibre reste
une prémisse mathématique précise --- le comptage uniforme de tous les
chemins compatibles --- et ne doit pas être confondue avec le successeur d'une seule
trajectoire. Le passage à la limite inverse de toutes les profondeurs conserve un
énoncé propre, formulé à la fin de la section.

\section{\(5\) niveaux de la notion de chemin}

Soit \(\Gamma_{\rm APP}\) la catégorie des chemins cardinaux du tore APP\@.
Pour chaque profondeur \(m\), soit \(\mathcal H_m\) l'ensemble fini des
histoires enrichies admissibles
\[
 h_m=(v_0,e_1,v_1,\ldots,e_m,v_m),
\]
où chaque flèche \(e_j\) transporte tous les champs déclarés par TPK\@. Si
\(n\geq m\), l'application
\[
 p_{n,m}:\mathcal H_n\longrightarrow\mathcal H_m
\]
élimine la queue postérieure à la profondeur \(m\). L'ombre positionnelle
\[
 \operatorname{sh}_m:\mathcal H_m\longrightarrow\Gamma_{\rm APP}
\]
oublie la fibre de mémoire et conserve le chemin APP sous-jacent.

\begin{definition}[Hiérarchie des extensions]
\label{def:jerarquia-extensiones}
Nous distinguons les objets suivants.
\begin{enumerate}
  \item Un \emph{chemin APP} est un morphisme de
  \(\Gamma_{\rm APP}\).
  \item Une \emph{route observable} est une lecture de ce morphisme qui
  conserve conjointement les feuillets de somme et de produit.
  \item Une \emph{extension TPK} est une élévation de la route observable en
  une histoire \(h_m\in\mathcal H_m\), avec ses données de résidu, de quotient,
  de signature, d'orientation, de frontière, de retenue et de registre.
  \item Une \emph{extension survivante} est une histoire finie qui admet
  des prolongations compatibles à tout horizon.
  \item Une \emph{section globale} est une famille compatible
  \(x=(h_m)_{m\geq0}\) d'extensions survivantes.
\end{enumerate}
En particulier, la route est l'ombre de l'extension ; ce n'est pas un objet extérieur
que l'extension choisirait ensuite.
\end{definition}

Pour \(h\in\mathcal H_m\), son cylindre de prolongations est
\[
 Z_m(h)=
 \{x=(h_n)_n\in\varprojlim_n\mathcal H_n:x_m=h\}.
\]
En particulier, toute famille de \(Z_m(h)\) satisfait en outre
\(p_{n,\ell}(h_n)=h_\ell\) pour \(n\geq\ell\) ; on n'admet pas de collections de
troncatures incompatibles.
Définissons également
\[
 \mathcal S_m^{(N)}=p_{N,m}(\mathcal H_N),
 \qquad
 \mathcal S_m^\infty=\bigcap_{N\geq m}\mathcal S_m^{(N)}.
\]

\begin{theorem}[Critère conditionnel d'extension survivante]
\label{thm:criterio-extension-superviviente}
Fixons \(h\in\mathcal H_m\) et, pour \(N\geq m\), écrivons
\[
 \mathcal P_N(h)=\{g\in\mathcal H_N:p_{N,m}(g)=h\}.
\]
Supposons que les troncatures sont cohérentes et que chaque
\(\mathcal P_N(h)\) est fini. Sont équivalentes :
\begin{enumerate}
  \item \(\mathcal P_N(h)\ne\varnothing\) pour tout \(N\geq m\), c'est-à-dire
  \(h\in\mathcal S_m^\infty\) ;
  \item l'arbre des prolongations finies de \(h\) est à ramification finie
  et possède un niveau non vide à toute profondeur ;
  \item \(Z_m(h)\ne\varnothing\).
\end{enumerate}
Par conséquent, dès lors que les troncatures survivantes finies sont non
vides, l'ensemble
\[
 \Omega_{\rm H}^{\rm surv}
 =\varprojlim_m\mathcal S_m^\infty
\]
est exactement l'espace des sections globales survivantes. Le théorème
n'affirme pas que ces troncatures sont non vides pour toute entrée APP : cette
propriété doit être démontrée au moyen du protocole de survie.
\end{theorem}

\begin{proof}
L'implication de la troisième condition vers les deux premières s'obtient par
troncature. Supposons maintenant la première condition et ordonnons les éléments de
\(\bigsqcup_{N\geq m}\mathcal P_N(h)\) par la relation de préfixe. Chaque niveau
est fini et non vide, et la cohérence des troncatures fait de cette union
un arbre à ramification finie avec des nœuds à toute profondeur ; cela donne la
deuxième condition. Réciproquement, à partir de la deuxième condition, le lemme de König
fournit une branche infinie. Ses troncatures forment une famille compatible
appartenant à \(Z_m(h)\), ce qui est la troisième condition. La même construction
identifie la limite inverse indiquée là où ses troncatures finies sont
non vides.
\end{proof}

Le théorème résout l'éventuelle discordance entre « prolongeable jusqu'à chaque
horizon séparément » et « possède une prolongation globale » : la finitude des
branches est l'hypothèse qui permet de choisir une unique chaîne compatible. Il
n'affirme pas que tous les chemins APP appartiennent à
\(\operatorname{sh}(\Omega_{\rm H}^{\rm surv})\). Cette couverture exige de démontrer
que les fermetures enrichies ne vident pas les fibres correspondantes.

\section{Retour de phase et relation de monodromie}

Le retour du quotient phase--longueur est formalisé séparément dans
\cref{def:odometro-9-12,prop:proyeccion-fase-longitud} :
\[
(b,r)\xrightarrow{\;9\;}(b+1,r),
\qquad
(b,r)\xrightarrow{\;108\;}(b,r).
\]
Ces égalités appartiennent à l'odomètre \(9\times12\). Elles ne constituent pas la
démonstration de \(F_{\rm obs}^{108}=I\), qui utilise l'état observable et apparaît plus
loin.

Soit \(\phi\) la coordonnée de phase à valeurs dans \(\mathbb Z/9\mathbb Z\).
Un retour visible satisfait
\[
 \phi(h_{k+9})=\phi(h_k).
\]
Le registre complet peut cependant satisfaire
\[
 B(h_{k+9})\ne B(h_k).
\]
Cette inégalité n'est pas un défaut de périodicité : c'est la mémoire du tour.
Une action ordinaire de \(C_9\) sur l'état complet imposerait que la
neuvième puissance soit l'identité et effacerait précisément cette information.

Pour typer le retour, soit \(H\) une hexade et
\(P_2(H)\) son ensemble de quinze paires. Soient \(\mathcal B_k\) les états de
registre qui survivent à la phase \(k\), soit
\(\mathcal T_k\subseteq\mathcal B_k\times\mathcal B_{k+1}\) la relation de
transport publiée et supposons donné un encodeur de fibre
\[
 c_k:\mathcal B_k\longrightarrow P_2(H).
\]
Ce n'est qu'après avoir déclaré ces encodeurs que l'on définit la relation observable
\[
 R_k=(c_k\times c_{k+1})(\mathcal T_k)
 \subseteq P_2(H)\times P_2(H).
\]
La monodromie relationnelle d'un tour est
\[
 \mathcal M_k
 =R_{k+8}\circ R_{k+7}\circ\cdots\circ R_k.
\]

\begin{proposition}[Classification du retour nonadique étant donné l'encodeur]
\label{prop:clasificacion-retorno-nonadico}
On distingue trois niveaux logiquement distincts.
\begin{enumerate}
  \item Sans démonstration de fonctionnalité, le retour est la relation
  \(\mathcal M_k\).
  \item Si chaque \(R_k\) est le graphe d'une application
  \(\mu_k:P_2(H)\to P_2(H)\), le transport est le produit biaisé
  \[
   T(k,p)=(k+1,\mu_k(p)),
  \]
  et
  \[
   T^9(k,p)=(k,M_k(p)),
   \qquad
   M_k=\mu_{k+8}\circ\cdots\circ\mu_k.
  \]
  \item Si les \(\mu_k\) sont bijectives, \(M_k\) est une holonomie
  automorphe. Le transport descend en une action ordinaire de
  \(C_9\) sur l'état complet si et seulement si \(M_k=I\) dans toutes les
  fibres.
\end{enumerate}
\end{proposition}

\begin{proof}
La première affirmation est la définition de l'image et de la composition relationnelle une
fois fixés \((c_k)_k\). Sous
fonctionnalité, neuf applications consécutives maintiennent la première
coordonnée fixe et composent leurs actions sur la seconde. La composition de
bijections est une bijection. Enfin, \(T^9=I\) équivaut exactement à
\(M_k=I\) pour chaque phase et chaque point de la fibre.
\end{proof}

\begin{corollary}[Non-retour ponctuel sous transport accumulé]
\label{pf84:E02-29}
Soient \(q:X\to V\), \(T:X\to X\) et \(x\in X\). Si
\[
 q(T^9x)=q(x),\qquad
 T^9x=\mathcal M_x\,x,\qquad
 \mathcal M_x\,x\ne x,
\]
alors la projection revient, mais pas l'état :
\[
 q(T^9x)=q(x),\qquad T^9x\ne x.
\]
L'hypothèse ponctuelle \(\mathcal M_x\,x\ne x\) est indispensable ; l'inégalité
globale \(\mathcal M_x\ne I\) n'exclut pas que \(x\) soit un point fixe.
\end{corollary}

\begin{proof}
La première égalité est une hypothèse. La seconde et la séparation ponctuelle donnent
\(T^9x=\mathcal M_x\,x\ne x\). L'énoncé serait réfuté par un cas qui
satisferait les trois prémisses et aurait \(T^9x=x\) ; il n'autorise pas à remplacer la
troisième prémisse par \(\mathcal M_x\ne I\).
\end{proof}

\begin{theorem}[Suffisance fonctionnelle du tuple enrichi]
\label{pf84:E02-30}
Soit
\[
 E:\mathscr X_{\rm TPK}^{\rm enr}\longrightarrow\mathscr S_E
\]
le tuple qui conserve le bloc visible, les blocs nouveaux et la retenue, l'état
d'élévation, les cylindres, la signature de frontière, l'orientation, l'observable, la phase et le
registre dodécaphasique. Soient
\[
 T:\mathscr X_{\rm TPK}^{\rm enr}\to
 \mathscr X_{\rm TPK}^{\rm enr},
 \qquad
 O:\mathscr X_{\rm TPK}^{\rm enr}\to Y.
\]
Lorsqu'il existe des applications
\[
 \widehat T:\mathscr S_E\to\mathscr X_{\rm TPK}^{\rm enr},
 \qquad
 \widehat O:\mathscr S_E\to Y
\]
tels que
\[
 T=\widehat T\circ E,\qquad O=\widehat O\circ E,
\]
l'égalité \(E(x)=E(y)\) implique \(T(x)=T(y)\), \(O(x)=O(y)\) et, par
déterminisme, les mêmes continuations et sorties ultérieures.
\end{theorem}

\begin{proof}
Par factorisation,
\[
 T(x)=\widehat T(E(x))=\widehat T(E(y))=T(y),
 \qquad
 O(x)=\widehat O(E(x))=\widehat O(E(y))=O(y).
\]
La première égalité permet d'itérer l'argument. C'est une condition
suffisante, non une équivalence non démontrée. Deux états ayant le même tuple et
des transitions ou des sorties différentes, l'accès à une coordonnée absente ou la
consultation d'une valeur cible sont des réfutateurs. La microfibre visible de
cardinal cinq de \cref{pf84:E02-22} contrôle uniquement la non-fidélité du
mot \(w_6\) ; elle ne démontre pas cette factorisation positive.
\end{proof}

La fenêtre finie publiée contient neuf retours de phase avec changement de
l'état enrichi et une sélection locale passant de trois continuations à une lors de
l'augmentation de deux horizons. Cela démontre le retour avec mémoire et réfute
l'identification du registre à un cycle sans mémoire. L'extraction globale des
neuf encodeurs \(c_k\) et, à travers eux, des relations
\(R_k\), à partir de toutes les histoires survivantes est une opération distincte.
La proposition précédente classe le retour une fois donnée l'application
registre--paires ; elle ne dérive pas cette application de l'état enrichi du TPK.

L'épisode concret de la neuvième phase, avec ses trois continuations, le
bloc \(B_{11}\), la frontière \((502,662,638)\) et les inégalités entières
de cylindres, est démontré une seule fois dans
\cref{thm:seleccion-novena-fase}. Ici, seule sa
conséquence typée est utilisée :
\[
\#\operatorname{Surv}_{1}(B_9)=3,\qquad
\#\operatorname{Surv}_{2}(B_9)=1.
\]

Les deux autres mots ne peuvent précéder la même queue ternaire qu'en changeant
l'extension profonde de Hensel. Ils partagent une partie du résidu visible,
mais non la généalogie de l'état.

\section{Le catalogue de 468 émissions comme produit de canaux}

Le catalogue fini part de l'ensemble
\(\Omega_{\rm seed}\) de \(104\,976\) présentations microscopiques. Soit
\(\mathcal U_6\) l'ensemble des émissions sextuples distinctes et soit
\[
 F:\Omega_{\rm seed}\twoheadrightarrow\mathcal U_6.
\]
Le recensement exact est
\[
 |\mathcal U_6|=468,
 \qquad
 \#\{U:|F^{-1}(U)|=144,432,1944\}=432,18,18.
\]
Ces comptages appartiennent à \(F\). Les comparer à des histoires TPK de
profondeur six nécessiterait une application explicite du domaine des histoires
vers \(\Omega_{\rm seed}\) ; aucun résultat de cette section ne présuppose cette
application ni n'identifie des présentations du catalogue à des états enrichis.
Le catalogue fini fournit, pour chaque émission, une signature additive
\(s\) et une signature multiplicative \(e\). L'application conjointe est une
bijection
\[
 \mathcal U_6\xrightarrow{\sim}
 \mathcal S_+\times\mathcal S_\times,
 \qquad
 |\mathcal S_+|=18,
 \quad
 |\mathcal S_\times|=26.
\]
Ainsi,
\[
 \boxed{468=18\cdot26}
\]
est une factorisation de l'objet et non seulement de son cardinal. La première
coordonnée conserve la manière dont l'émission additionne ; la seconde, la manière dont elle accumule le
produit. Cette séparation est la forme finie explicite de
l'incompatibilité somme--produit : les deux lectures coexistent dans l'état, mais
chaque projection oublie la signature complémentaire.

La première coordonnée admet une paramétrisation supplémentaire qui sera utile pour
la comparer à la frontière nonadique. Si un germe additif a pour coordonnées
\((i,j)\in(\mathbb Z/9\mathbb Z)^2\) et pour orientation
\(\varepsilon=+1\) pour les directions est--sud et \(-1\) pour
nord--ouest, sa signature complète est déterminée bijectivement par
\[
 \bigl(i+j\bmod9,\varepsilon\bigr)
 \in(\mathbb Z/9\mathbb Z)\times\{\pm1\}.
\]
Par conséquent,
\[
 \boxed{\mathcal S_+\simeq
 (\mathbb Z/9\mathbb Z)\times\{\pm1\}.}
\]
La coordonnée produit ne possède pas une uniformité analogue : ses \(26\) signatures
ont des fibres de germes de tailles \(8\), \(24\) et \(108\), avec l'histogramme
\(24\times8+1\times24+1\times108=324\). Cette asymétrie est structurelle. Le
facteur \(18\) contient déjà la phase et l'orientation ; toute descente ultérieure
du facteur produit vers une classe de paires doit conserver une mémoire supplémentaire ou
accepter des fibres non uniformes.

Définissons sur \(\ell^2(\mathcal S_+\times\mathcal S_\times)\) les
espérances conditionnelles
\[
\begin{aligned}
 (E_+f)(s,e)&=\frac1{18}\sum_{s'\in\mathcal S_+}f(s',e),\\
 (E_\times f)(s,e)&=\frac1{26}\sum_{e'\in\mathcal S_\times}f(s,e').
\end{aligned}
\]
Le premier opérateur renouvelle la signature additive et conserve la multiplicative ;
le second réalise l'opération duale. Le troisième canal retient les deux.

\section{La trajectoire unique et le transfert de chemins}

À la différence du retour de l'odomètre défini dans
\cref{def:odometro-9-12,prop:proyeccion-fase-longitud}, la proposition
suivante démontre le retour de l'opérateur observable \(F_{\rm obs}\).

La loi observable publiée met à jour un seul curseur à chaque pas. Aux
phases \(+1\), le curseur additif avance ; aux phases \(-1\), le
multiplicatif ; aux phases \(0\), les deux restent immobiles. Après les
pas \(27,54,81,108\), les deux orientations sont inversées. Notons cette
évolution déterministe \(F_{\rm obs}\).

\begin{proposition}[Retour de la chronologie observable]
\label{prop:retorno-fobs-identidad}
Pour tout choix des deux positions et orientations initiales, les
positions et les orientations sont restaurées après \(54\) pas. En
restaurant également la phase, on obtient
\[
 F_{\rm obs}^{108}=I.
\]
Par conséquent, le retour de Poincaré de \(108\) pas induit l'identité sur
\(\mathcal U_6\) et ne sélectionne pas une mesure stationnaire unique.
\end{proposition}

\begin{proof}
Dans chaque bloc de \(27\) pas, chaque curseur effectue exactement neuf
avancements dans une orientation fixe. Comme l'espace des positions est
\(X_9=(\mathbb Z/9\mathbb Z)^2\), son déplacement est \(9d=0\). À la fin
du bloc, \(d\) est inversé. Deux blocs restaurent la position et l'orientation ;
quatre restaurent aussi l'indice dans \(\Theta_{108}\). Le retour induit
est donc l'identité. Toute probabilité est stationnaire pour
l'identité, de sorte qu'aucune n'est sélectionnée par cette chronologie.
\end{proof}

Le problème apparaît même avant le retour si l'on tente d'oublier la
mémoire de route. La signature produit
\[
 e_0=(5,5,5,5,5,5)
\]
possède vingt-quatre représentants. Après un avancement dans leur propre orientation,
ces représentants se séparent, huit par huit, entre
\[
 (5,5,5,7,1,7),\qquad
 (5,5,5,7,7,1),\qquad
 (5,5,5,1,7,7).
\]
Ainsi, l'avancement produit ne descend pas en une application des vingt-six
signatures : le représentant qui avait été oublié redevient pertinent. Ce
témoin explique pourquoi une route individuelle ne peut pas devenir, par simple
projection, le renouvellement d'une fibre complète.

La lecture APP est différente. Au lieu de choisir un représentant, elle conserve toutes
les continuations compatibles. Pour la formaliser dans le catalogue fini,
considérons les sous-algèbres de fonctions qui ne dépendent que de l'une des deux
signatures. Par rapport au comptage uniforme, \(E_+\) et \(E_\times\) sont les
espérances conditionnelles qui oublient exactement la coordonnée active. Ce sont les
seuls projecteurs de Markov qui satisfont simultanément :
\begin{enumerate}
  \item la préservation de la coordonnée inactive ;
  \item l'oubli complet de la coordonnée active ;
  \item la préservation de l'état de comptage uniforme.
\end{enumerate}
En effet, une ligne qui a oublié la coordonnée de départ doit être
constante sur les \(18\), respectivement \(26\), destinations compatibles ; la
normalisation impose les valeurs \(1/18\) et \(1/26\). C'est la canonicité
relative du renouvellement : on ne postule pas de préférence entre des branches qu'APP
déclare également possibles.

Soit \(C_9=\mathbb Z/9\mathbb Z\), avec le mot de phase
\[
 \tau=(+1,-1,0,+1,-1,0,+1,-1,0).
\]
Écrivons
\[
 P_{+1}=E_+,
 \qquad P_{-1}=E_\times,
 \qquad P_0=I.
\]

\begin{definition}[Transfert TPK biaisé par la phase]
\label{def:transferencia-tpk-fase}
Sur \(\mathcal U_6\times C_9\), nous définissons
\[
 \mathcal K_{\rm ph}\bigl((u,r),(v,r+1)\bigr)
 =P_{\tau(r)}(u,v),
\]
et nous annulons les autres entrées. Cette définition est relative au lecteur de
tous les chemins compatibles : une phase active renouvelle uniformément la fibre
que cette phase lit et la phase neutre retient l'émission.
\end{definition}

\begin{theorem}[Retour nonadique uniforme]
\label{thm:retorno-nonadico-uniforme}
Le noyau \(\mathcal K_{\rm ph}\) est doublement stochastique et irréductible,
a une période exactement égale à neuf et possède un unique état stationnaire,
\[
 \pi_{\rm ph}(u,r)=\frac1{9\cdot468}.
\]
Dans chaque fibre de phase, son retour de neuf pas est
\[
 \mathcal K_{\rm ph}^{\,9}\big|_r=E_+E_\times,
 \qquad
 (E_+E_\times)(u,v)=\frac1{468}.
\]
Si la phase initiale est distribuée uniformément, l'ombre d'un pas sur
les émissions est
\[
 \overline K_{\rm ph}=\frac13(I+E_++E_\times).
\]
\end{theorem}

\begin{proof}
Chaque \(P_{\tau(r)}\) est doublement stochastique et le déplacement
\(r\mapsto r+1\) est une permutation ; leur produit biaisé conserve les deux sommes.
Les trois opérateurs sont idempotents et commutent. Comme chacun apparaît trois
fois dans le mot nonadique,
\[
 \prod_{r\in C_9}P_{\tau(r)}
 =E_+^3E_\times^3I^3=E_+E_\times.
\]
La première espérance oublie la signature additive et la seconde la multiplicative,
de sorte que leur composition attribue à chacune des \(18\cdot26=468\)
destinations la probabilité \(1/468\). Ainsi, en neuf pas, deux
états quelconques d'une même phase communiquent ; les pas restants font communiquer les
phases, et le noyau est irréductible. Tout retour doit avoir une longueur multiple
de neuf, tandis que la diagonale de \(\mathcal K_{\rm ph}^9\) est positive ;
la période est neuf. La double stochasticité et l'irréductibilité donnent
l'unique état uniforme. Enfin, moyenner les neuf phases compte trois
copies de chaque canal et produit la dernière formule.
\end{proof}

Oublier la phase n'est pas un quotient de Markov fort : à partir d'une même émission,
deux phases distinctes appliquent \(I\), \(E_+\) ou \(E_\times\), qui ont des lignes
distinctes. L'opérateur \(\overline K_{\rm ph}\) est une ombre moyennée sur la
phase stationnaire, non le successeur chronologique d'une émission sans phase. Cette
distinction conserve à la fois la dynamique et le sens de la mesure.

\begin{definition}[Noyau moyenné induit par le transfert]
\label{def:nucleo-tritico-refresco}
Nous appelons noyau moyenné du catalogue l'ombre de phase
\[
 \boxed{
 K_{\rm cat}=\frac13(I+E_++E_\times).
 }
\]
Les poids sont induits par le calendrier TPK, qui contient trois occurrences
de chaque signe. Les opérateurs de renouvellement sont induits relativement à la
normalisation uniforme de toutes les continuations compatibles de la fibre
active. Aucune des deux affirmations n'identifie \(K_{\rm cat}\) à
\(F_{\rm obs}\).
\end{definition}

\begin{proposition}[Famille convexe de noyaux stochastiques à mesure uniforme commune]
\label{prop:familia-refrescos-abc}
Pour \(a,b,c\geq0\), \(a+b+c=1\), la famille
\[
 K_{a,b,c}=aI+bE_++cE_\times
\]
est formée de noyaux symétriques et doublement stochastiques. Si \(b,c>0\),
ils sont irréductibles et, ayant une diagonale positive, apériodiques. Leur spectre est
\[
 \operatorname{Spec}(K_{a,b,c})=
 \{1^{[1]},(a+c)^{[17]},(a+b)^{[25]},a^{[425]}\}.
\]
En particulier, il existe un continuum de dynamiques ayant la même factorisation et la
même mesure uniforme. Le choix
\(K_{\rm cat}=K_{1/3,1/3,1/3}\) n'est pas caractérisé par le seul recensement. Il est
sélectionné conjointement par le mot de phase TPK et le transfert
normalisé de \cref{def:transferencia-tpk-fase}.
\end{proposition}

\begin{proof}
Les trois termes sont des noyaux symétriques et stochastiques ; par convexité, il en va
de même pour \(K_{a,b,c}\), et la symétrie transforme la stochasticité par lignes
en stochasticité par colonnes. Si \(b,c>0\), deux pas permettent de renouveler
successivement les deux coordonnées et de faire communiquer toute paire d'états. De plus,
\[
 K_{a,b,c}((s,e),(s,e))=a+\frac b{18}+\frac c{26}>0,
\]
de sorte que le noyau est apériodique. Dans la décomposition orthogonale
\[
 \mathbf1\oplus\ell_0^2(\mathcal S_+)
 \oplus\ell_0^2(\mathcal S_\times)
 \oplus\bigl(\ell_0^2(\mathcal S_+)\otimes
              \ell_0^2(\mathcal S_\times)\bigr),
\]
les paires \((E_+,E_\times)\) agissent, respectivement, comme
\((1,1),(0,1),(1,0),(0,0)\). On obtient ainsi les quatre valeurs propres et les
multiplicités \(1,17,25,425\). Pour \(b,c>0\), l'irréductibilité rend unique
la mesure stationnaire uniforme, ce qui montre la non-unicité même sous
cette exigence.
\end{proof}

\begin{theorem}[Spectre exact du renouvellement équipondéré]
\label{thm:dinamica-cociente-468}
Le noyau \(K_{\rm cat}\) est symétrique, doublement stochastique,
irréductible et apériodique. Son spectre, avec multiplicités, est
\[
 \boxed{
 \operatorname{Spec}(K_{\rm cat})
 =\left\{1^{[1]},
 \left(\frac23\right)^{[42]},
 \left(\frac13\right)^{[425]}\right\}.
 }
\]
Par conséquent, son écart spectral est (1/3) et son unique état
stationnaire est
\[
 \boxed{
 \tau_{468}(f)=\frac1{468}\sum_{U\in\mathcal U_6}f(U).
 }
\]
\end{theorem}

\begin{proof}
Les opérateurs \(E_+\) et \(E_\times\) sont des projecteurs orthogonaux de
moyennage et commutent. Les lignes et les colonnes de chacun ont pour somme un ; il en va de même
pour leur moyenne. De \((s,e)\), on atteint \((s',e')\) en deux
renouvellements, et la contribution de \(I\) fournit une boucle. Cela
démontre l'irréductibilité et l'apériodicité.

L'espace se décompose orthogonalement en quatre secteurs :
\[
 \mathbf1,
 \qquad
 \ell_0^2(\mathcal S_+),
 \qquad
 \ell_0^2(\mathcal S_\times),
 \qquad
 \ell_0^2(\mathcal S_+)\otimes
 \ell_0^2(\mathcal S_\times).
\]
Leurs dimensions sont \(1,17,25,17\cdot25\). Dans le secteur constant, les
trois termes agissent comme l'identité ; dans chaque secteur d'une seule coordonnée,
ils agissent comme (I+I+0) ; dans le secteur d'interaction, comme (I+0+0).
Diviser par trois donne les valeurs propres et les multiplicités annoncées. Un noyau
fini irréductible et apériodique possède un unique état stationnaire ; comme il est
doublement stochastique, celui-ci est uniforme.
\end{proof}

Le théorème est exact relativement à la factorisation
\(\mathcal U_6\simeq\mathcal S_+\times\mathcal S_\times\), au calendrier
TPK et au comptage uniforme des continuations compatibles de chaque fibre. La
famille \(K_{a,b,c}\) démontre que la dynamique ne se déduit pas du seul
recensement. \cref{prop:retorno-fobs-identidad} distingue le transfert de
tous les chemins de la chronologie de l'un d'entre eux.

\section{Élévation compensée construite à partir du quotient}

Écrivons \(m(U)=|F^{-1}(U)|\). À partir du noyau déjà choisi dans le quotient,
nous définissons le noyau de présentations par
\[
 \boxed{
 K_{\rm seed}(x,y)
 =\frac{K_{\rm cat}(F(x),F(y))}{m(F(y))}
 }
\]
Ce noyau est stochastique. En sommant sur une fibre de destination, on obtient
\[
 \sum_{F(y)=V}K_{\rm seed}(x,y)
 =K_{\rm cat}(F(x),V),
\]
indépendamment du représentant \(x\). Cette identité montre que
l'élévation construite se projette exactement sur \(K_{\rm cat}\). Elle ne démontre pas
qu'une dynamique microscopique préexistante de l'état enrichi du TPK descende
au catalogue.

\begin{theorem}[Élévation réversible construite à partir du quotient]
\label{thm:elevacion-reversible-cociente}
La probabilité
\[
 \widetilde\tau(x)
 =\frac1{468\,m(F(x))}
\]
est stationnaire et réversible pour \(K_{\rm seed}\), et
\(F_\#\widetilde\tau=\tau_{468}\). Si
\[
 J_F\delta_U
 =m(U)^{-1/2}\sum_{F(x)=U}\delta_x,
 \qquad
 D=\operatorname{diag}(m(U)),
\]
alors
\[
 K_{\rm seed}J_F
 =J_F\bigl(D^{1/2}K_{\rm cat}D^{-1/2}\bigr).
\]
En particulier, l'entrelacement hilbertien sans conjugaison par \(D\)
a lieu si et seulement si \(K_{\rm cat}\) commute avec les multiplicités des
fibres.
\end{theorem}

\begin{proof}
La normalisation découle du fait que chacune des 468 fibres reçoit une masse
totale (1/468). En utilisant la symétrie de \(K_{\rm cat}\), pour
\(F(x)=U\) et \(F(y)=V\), on a
\[
 \widetilde\tau(x)K_{\rm seed}(x,y)
 =\frac{K_{\rm cat}(U,V)}{468m(U)m(V)}
 =\widetilde\tau(y)K_{\rm seed}(y,x).
\]
Cela démontre l'équilibre détaillé, la stationnarité et l'identité de la
mesure image.
Appliquer les deux membres de la dernière identité à un vecteur de base
\(\delta_V\) donne, sur la fibre \(U\), le coefficient
\(K_{\rm cat}(U,V)/\sqrt{m(V)}\) ; c'est exactement le coefficient de
l'opérateur conjugué présenté.
\end{proof}

\begin{theorem}[Élévation de phase fidèle à l'immobilité]
\label{thm:elevacion-fase-semillas}
Soit \(X=\Omega_{\rm seed}\), \(m(u)=|F^{-1}(u)|\), et définissons
\[
\begin{aligned}
 L_0(x,y)&=\mathbf1_{\{x=y\}},\\
 L_+(x,y)&=\frac{E_+(F(x),F(y))}{m(F(y))},\\
 L_\times(x,y)&=\frac{E_\times(F(x),F(y))}{m(F(y))}.
\end{aligned}
\]
Sur \(X\times C_9\), soit
\[
 \widehat{\mathcal K}_{\rm ph}\bigl((x,r),(y,r+1)\bigr)
 =L_{\tau(r)}(x,y).
\]
Alors \((F,\mathrm{id}_{C_9})\) satisfait la propriété
d'\emph{agrégabilité forte} (\emph{strong lumpability}) de
\(\widehat{\mathcal K}_{\rm ph}\) vers \(\mathcal K_{\rm ph}\). Le noyau
microscopique est irréductible, a pour période neuf et son unique état
stationnaire est
\[
 \widetilde\pi_{\rm ph}(x,r)
 =\frac1{9\cdot468\,m(F(x))}.
\]
Son retour nonadique satisfait
\[
 \widehat{\mathcal K}_{\rm ph}^{\,9}
 \bigl((x,r),(y,r)\bigr)
 =\frac1{468\,m(F(y))}.
\]
\end{theorem}

\begin{proof}
La somme de \(L_+\), respectivement \(L_\times\), sur une fibre de destination
est \(E_+\), respectivement \(E_\times\), et \(L_0\) se projette sur l'identité.
Cela démontre l'agrégabilité à chaque phase. Les facteurs \(m(F(y))\) s'annulent
lors de la composition : on sélectionne d'abord uniformément la signature active, puis un
représentant de l'émission obtenue. Ainsi, \(L_+\) et \(L_\times\) sont des
projecteurs commutants et
\[
 L_+L_\times(x,y)=\frac1{468\,m(F(y))}.
\]
La positivité de cette entrée fait communiquer tous les représentants en neuf
pas. La phase impose que les longueurs de retour soient des multiples de neuf,
et l'entrée diagonale du retour est positive. La formule stationnaire se
vérifie en sommant d'abord les \(m(u)\) représentants de chaque émission ; chacune
des \(9\cdot468\) classes phase--émission reçoit la même masse.
\end{proof}

Cette élévation améliore la lecture du noyau moyenné
\(K_{\rm seed}\) : dans la phase neutre, elle conserve le représentant microscopique,
au lieu de remélanger silencieusement sa fibre. En moyennant et en oubliant la phase,
les deux constructions ont la même ombre \(K_{\rm cat}\), mais pas la même
dynamique cachée. La différence est nécessaire pour respecter le fait que le zéro de TRIT
est une mémoire d'immobilité.

\section{Congruences minimales stables par avancement, demi-tour et 4e de tour : \(28\) et \(162\) classes de mémoire}

Le témoin \(e_0=(5,5,5,5,5,5)\) montre que la signature produit est trop
grossière pour transporter une route. Il est possible de calculer exactement la mémoire
minimale manquante. Soit
\[
 \mathcal X_\times=X_9\times\mathscr D_{\rm card}
\]
l'espace des \(324\) germes produit. Notons \(P\) l'avancement d'une
case dans l'orientation stockée, \(H\) le demi-tour de cette
orientation et \(Q\) le quart de tour
\[
 N\longmapsto E\longmapsto S\longmapsto O\longmapsto N.
\]

\begin{theorem}[Congruences d'avancement et de rotation]
\label{thm:congruencias-avance-giro}
La congruence minimale qui raffine les vingt-six signatures produit et est stable
par \(P\) et \(H\) possède \(28\) classes : vingt-sept sont de taille huit et une
est de taille \(108\). L'unique scission nouvelle est
\[
 (5,5,5,5,5,5)\longrightarrow
 \mathcal L_0\sqcup\mathcal L_1\sqcup\mathcal L_2,
 \qquad |\mathcal L_j|=8.
\]
Ses orbites sous \(P,H\) ont pour tailles
\[
 9,\quad9,\quad9,\quad1.
\]
En combinant ce raffinement avec les dix-huit signatures additives, on obtient
\[
 18\cdot28=504=468+36.
\]

La congruence minimale stable par \(P\) et \(Q\) possède \(162\) classes de deux
germes. Chaque classe est une orbite de l'involution
\[
 J(i,j,d)=\bigl(7-i,\,7-j,\,\overline d\bigr)\pmod9.
\]
Les opérateurs induits par \(P,Q\) agissent transitivement sur les
\(162\) classes.
\end{theorem}

\begin{proof}
On part de la partition par les vingt-six signatures et l'on raffine
itérativement une classe dès que deux de ses éléments ont des images dans des
classes distinctes par l'un des générateurs. Le processus se termine car
\(\mathcal X_\times\) est fini. Le recensement exhaustif produit, pour
\((P,H)\), l'histogramme \(27\times8+1\times108=324\), les quatre orbites
indiquées et trois classes au-dessus de la fibre exceptionnelle de taille vingt-quatre.
Pour \((P,Q)\), il produit \(162\times2=324\). La vérification directe démontre
que \(J\) préserve chaque signature, commute avec les deux générateurs et apparie
exactement les deux éléments de chaque classe ; le graphe de classes engendré par
\(P,Q\) est connexe.
\end{proof}

L'excès \(36\) et l'objet \(R_{36}\) sont de types distincts. La frontière
de profondeur trente-six est
\[
 R_{36}=
 \begin{pmatrix}
 222220\\021222\\102011
 \end{pmatrix},
\]
c'est-à-dire trois blocs ordonnés de six trits, un par canal. Ce n'est pas un
ensemble de trente-six états. De son côté, le terme
\(504-468=36\) compte une mémoire supplémentaire de feuillet : le demi-tour qui rend
fonctionnel le transport produit ajoute exactement deux feuillets nouveaux à une
seule signature et, en les répétant sur les dix-huit signatures additives, produit
\(18(3-1)=36\). La coïncidence de l'entier (36) exprime deux graduations
du transport, profondeur de frontière et multiplicité de feuillet, mais
n'autorise pas à identifier leurs objets.

Le quart de tour \(Q\) n'était pas non plus caché dans \(F_{\rm obs}\). La
chronologie publiée contient des avancements et des demi-tours programmés ; incorporer
\(Q\) est une extension explicite de l'alphabet des routes. Une fois incorporé,
une marche symétrique paresseuse dans \(P^{\pm1},Q^{\pm1}\) est irréductible et
apériodique sur les \(162\) classes, et sa mesure uniforme est unique. C'est la
forme mathématique exacte de « compter les rotations » : la rotation est une coordonnée de
transport et non une correction décimale ultérieure.

\section{Cocycles de mode et d'orientation}
\label{sec:cociclos-modo-orientacion}

La même règle permet de définir les compteurs opératoires sans introduire la
mécanique dimensionnelle. Ajoutons à l'état le dernier mode actif
\(m\in\{\Sigma,\Pi\}\). Une signature \(+1\) fixe \(m=\Sigma\), une signature
\(-1\) fixe \(m=\Pi\), et une signature \(0\) conserve le mode précédent. Pour une
flèche \(a\), définissons le cocycle signé
\[
 \omega_{\rm sw}(a)=
 \begin{cases}
  +1,&\Pi\longrightarrow\Sigma,\\
  -1,&\Sigma\longrightarrow\Pi,\\
  0,&\text{le mode est conservé}.
 \end{cases}
\]
Sa variation positive est \(s(a)=|\omega_{\rm sw}(a)|\). Définissons en outre
\(g(a)=1\) lorsque la paire ordonnée des orientations cardinales change et
\(g(a)=0\) lorsqu'elle est conservée. En maintenant aux extrémités le mode et
l'orientation terminale, les sommes
\[
 \Omega_{\rm sw}(\gamma)=\sum_{a\subset\gamma}\omega_{\rm sw}(a),
 \qquad
 S(\gamma)=\sum_{a\subset\gamma}s(a),
 \qquad
 G(\gamma)=\sum_{a\subset\gamma}g(a)
\]
sont additives par concaténation.

Dans un tour nonadique fermé, la suite des modes est
\[
 (\Sigma,\Pi,\Pi,\Sigma,\Pi,\Pi,\Sigma,\Pi,\Pi).
\]
Par conséquent,
\[
 \Omega_{\rm sw}(C_9)=0,\qquad S(C_9)=6:
\]
il y a trois changements \(\Pi\to\Sigma\), trois changements
\(\Sigma\to\Pi\) et trois maintiens. Dans le cycle de \(108\) pas,
\[
 \#(+1)=\#(-1)=\#(0)=36,\qquad
 \Omega_{\rm sw}=0,\qquad S=72,\qquad G=4.
\]
Le dernier nombre compte les quatre inversions simultanées programmées dans
\(27,54,81,108\). Ces cocycles proviennent directement du calendrier TPK\@.

\subsection{Quotient des modes et correspondance aréale--volumétrique}
\label{sec:cociente-modos-fav}

Le signe du calendrier, l'étiquette de mode et le feuillet géométrique sont trois
objets typés. En particulier, on n'identifie pas le signe de phase à
l'observable angulaire sur les chiffres. Pour la correspondance géométrique, nous fixons le
lecteur HMT typé
\[
 \mathfrak g_{\rm av}(\Sigma)=\mathscr H_{\rm ar},
 \qquad
 \mathfrak g_{\rm av}(\Pi)=\mathscr H_{\rm vol}.
\]
La première attribution exprime la lecture additive ou de bord ; la seconde, la
lecture multiplicative ou d'intérieur. L'application \(\mathfrak g_{\rm av}\) fait
partie du typage des feuillets. Le résultat suivant démontre qu'une fois
ce typage fixé, la dynamique TPK impose l'incidence entre eux.

\begin{theorem}[Descente nécessaire de l'incidence des feuillets]
\label{thm:descenso-necesario-fav}
Soit \(m_r\) le mode obtenu après application de la phase \(r\) du bloc
primitif \((+1,-1,0)\), avec la règle
\[
 +1:\ m\mapsto\Sigma,\qquad
 -1:\ m\mapsto\Pi,\qquad
 0:\ m\mapsto m.
\]
L'orbite cyclique des modes est \((\Sigma,\Pi,\Pi)\). Si
\[
 N_3(i,j)
 :=
 \#\{r\in\mathbb Z/3\mathbb Z:(m_r,m_{r+1})=(i,j)\},
\]
alors, dans l'ordre \((\Sigma,\Pi)\),
\[
 \boxed{
 N_3=
 \begin{pmatrix}0&1\\1&1\end{pmatrix}.}
\]
Par répétition du calendrier,
\[
 N_9=3N_3,\qquad N_{27}=9N_3,\qquad N_{108}=36N_3.
\]
Après application de \(\mathfrak g_{\rm av}\), la matrice primitive
d'incidence est, dans l'ordre
\((\mathscr H_{\rm ar},\mathscr H_{\rm vol})\),
\[
 \boxed{
 F_{\rm av}=
 \begin{pmatrix}0&1\\1&1\end{pmatrix}.}
\]
Ainsi, les deux transitions croisées, l'unique autorétention volumétrique,
l'absence d'autorétention aréale et les multiplicités primitives unitaires
sont des conséquences du quotient des modes de TPK ; ce ne sont pas quatre postulats
ajoutés.
\end{theorem}

\begin{proof}
Les trois couples consécutifs, avec fermeture cyclique, sont
\[
 \Sigma\longrightarrow\Pi,\qquad
 \Pi\longrightarrow\Pi,\qquad
 \Pi\longrightarrow\Sigma.
\]
Chacun apparaît une fois et le couple
\(\Sigma\to\Sigma\) n'apparaît pas. Cela détermine les quatre entrées de
\(N_3\). Les calendriers de longueurs \(9,27,108\) contiennent,
respectivement, \(3,9,36\) copies du bloc primitif, de sorte que leurs
matrices d'incidence sont les multiples indiqués. Le plus grand commun diviseur
des entrées non nulles de chaque matrice longue élimine uniquement cette
répétition globale et restitue \(N_3\). Enfin,
\(\mathfrak g_{\rm av}\) transporte \(\Sigma,\Pi\) vers les feuillets aréal et
volumétrique sans modifier l'incidence.
\end{proof}

La réalisation hyperbolique réelle qui utilise cette matrice et prouve son
invariance lorentzienne est formulée dans
\cref{sec:tres-generadores-concretos-trit} ; on y utilise \(F_{\rm av}\)
comme antécédent, sans répéter sa dérivation.

Ce théorème doit être distingué d'une affirmation de Markov plus forte.
L'oubli de la phase n'est pas un quotient de Markov fort de
\(\mathcal K_{\rm ph}\) : le mode visible ne détermine pas lequel des trois
opérateurs \(I,E_+,E_\times\) agit à l'itération élémentaire suivante. Ce qui descend
sans hypothèse supplémentaire est le multigraphe des arêtes d'un bloc TPK\@.

Il existe cependant une complétion à mémoire un déterminée de manière
unique par cette descente. Soit \(X_{\rm av}\) le plus petit sous-décalage de type fini
sur
\(\{\mathscr H_{\rm ar},\mathscr H_{\rm vol}\}\) qui contient tous les
mots obtenus en oubliant la phase et qui ne conserve que des couples consécutifs.
Tout sous-décalage à un pas ayant cette propriété doit admettre exactement les trois
couples observés ; le plus petit d'entre eux exclut l'unique couple non observé. Par
conséquent, sa matrice d'adjacence est \(F_{\rm av}\).

\begin{corollary}[Mesure de Parry de l'enveloppe à mémoire un]
\label{cor:parry-envolvente-fav}
Le sous-décalage \(X_{\rm av}\) est primitif, a pour entropie topologique
\(\log\varphi\) et possède une unique mesure d'entropie maximale. Ses poids
stationnaires de sommet satisfont
\[
 \boxed{
 \frac{\mu_{\rm ar}}{\mu_{\rm vol}}=\varphi^{-2}.}
\]
\end{corollary}

\begin{proof}
La racine de Perron de \(F_{\rm av}\) est \(\varphi\), et un vecteur de Perron
positif est \((1,\varphi)\). Comme la matrice est symétrique, les poids de sommet
de Parry sont proportionnels au produit des vecteurs gauche et droit,
c'est-à-dire à \((1,\varphi^2)\).
\end{proof}

La mesure de Parry appartient à l'enveloppe des chemins à mémoire un. Elle n'est pas
l'image de \(\tau_{468}\), ni la fréquence temporelle de l'orbite périodique
des modes. Cette dernière est
\[
 \nu_{\rm cyc}(\mathscr H_{\rm ar})=\frac13,\qquad
 \nu_{\rm cyc}(\mathscr H_{\rm vol})=\frac23,
\]
tandis que \(\tau_{468}\) est uniforme sur les émissions et que la mesure de Parry
pondère des histoires de renouvellement de longueur arbitraire. Les trois mesures
répondent à des questions distinctes. Le rapport irrationnel
\(\varphi^{-2}\) est ainsi déduit canoniquement de l'enveloppe des chemins,
sans l'attribuer à un quotient de cardinaux finis.

\subsubsection{Du comptage chronologique au retour marqué par \texorpdfstring{\(\pi\)}{π}}
\label{sec:retorno-marcado-pi-phi2}

Les deux formules
\[
 \left(\frac13,\frac23\right)
 \qquad\text{et}\qquad
 \frac{\mu_{\rm ar}}{\mu_{\rm vol}}=\varphi^{-2}
\]
ne sont pas des approximations l'une de l'autre. La première compte les trois instants du
calendrier primitif \((\Sigma,\Pi,\Pi)\) ; la seconde pondère toutes les
histoires admissibles de l'enveloppe à mémoire un. C'est pourquoi la fréquence
chronologique est rationnelle et le rapport de Parry est irrationnel. Dans ce qui suit,
nous écrivons
\[
 \frac{\pi_{\rm HMT}}{\varphi^2}
 =\pi_{\rm HMT}\varphi^{-2};
\]
il ne s'agit pas de diviser par \(\varphi^{-2}\).

L'apparition de \(\varphi^{-2}\) se voit directement dans la dynamique
linéaire. Si \(F_n\) désigne le \(n\)-ième nombre de Fibonacci,
\[
 F_{\rm av}^{\,n}
 =
 \begin{pmatrix}
  F_{n-1}&F_n\\
  F_n&F_{n+1}
 \end{pmatrix}
 \qquad(n\geq1).
\]
Les valeurs propres de \(F_{\rm av}\) sont
\(\varphi\) et \(-\varphi^{-1}\). La branche contractée inverse l'orientation en
une itération. Lors du passage à la mémoire quadratique, c'est-à-dire à
\(\operatorname{Sym}^2(\mathbb R^2)\), cette inversion est enregistrée avant la
disparition du signe. Dans la base \((x^2,2xy,y^2)\), en représentant par
lignes les images des vecteurs de base, l'opérateur induit est
\[
 \operatorname{Sym}^2(F_{\rm av})=
 \begin{pmatrix}
  0&0&1\\
  0&1&2\\
  1&1&1
 \end{pmatrix},
 \qquad
 \chi(t)=(t+1)(t^2-3t+1).
\]
Ses trois valeurs propres sont
\[
 \varphi^2,\qquad -1,\qquad \varphi^{-2}.
\]
Ainsi, \(\varphi^{-2}\) est le mode stable positif de la mémoire du second
ordre. Cette étape explique pourquoi le nombre d'or apparaît ici comme loi de
retour et non comme moyenne statistique ajoutée.

La clôture \(\pi\) appartient à un autre niveau du même état HMT : c'est la sortie
coinductive du lecteur régional déjà construit, non une nouvelle valeur propre de
\(F_{\rm av}\). Soient
\[
 P_{\rm ar}=
 \begin{pmatrix}1&0\\0&0\end{pmatrix},
 \qquad
 P_{\rm vol}=
 \begin{pmatrix}0&0\\0&1\end{pmatrix}.
\]
Parmi les opérateurs positifs qui préservent les deux feuillets, il en existe un et un
seul dont la normalisation volumétrique vaut un et dont la marque aréale est la clôture
engendrée \(\pi_{\rm HMT}\) :
\[
 D_\pi=\pi_{\rm HMT}P_{\rm ar}+P_{\rm vol}.
\]
L'unicité est relative à ces trois clauses explicites : positivité,
diagonalité par rapport aux feuillets et les deux normalisations. Avec
\(e_{\rm ar}=(1,0)^T\) et \(e_{\rm vol}=(0,1)^T\), le retour marqué de
longueur \(n\) est
\[
 \Theta_n
 :=
 \frac{\langle e_{\rm ar},
 D_\pi F_{\rm av}^{\,n}e_{\rm ar}\rangle}
 {\langle e_{\rm vol},
 F_{\rm av}^{\,n}e_{\rm vol}\rangle}
 =
 \pi_{\rm HMT}\frac{F_{n-1}}{F_{n+1}}.
\]
Par conséquent,
\[
 \boxed{\displaystyle
 \lim_{n\to\infty}\Theta_n
 =\frac{\pi_{\rm HMT}}{\varphi^2}.}
\]
La marque \(\pi_{\rm HMT}\) est appliquée une seule fois, au retour de frontière ;
elle n'est pas réinjectée à chaque transition. La combinatoire finie produit le mode
stable \(\varphi^{-2}\), le lecteur coinductif produit la clôture
\(\pi_{\rm HMT}\), et \(D_\pi\) compose les deux données sans inverser leur ordre
causal.

La longueur native \(108\) offre un témoin exact de cette convergence :
\[
 (F_{\rm av}^{108})_{\rm ar,ar}
 =F_{107}=10284720757613717413913,
\]
\[
 (F_{\rm av}^{108})_{\rm vol,vol}
 =F_{109}=26925748508234281076009.
\]
En conséquence,
\[
 \left|
 \frac{F_{107}}{F_{109}}-\varphi^{-2}
 \right|
 =
 6.1684979916614407936\ldots\times10^{-46},
\]
et, après l'unique marque de clôture,
\[
 \left|
 \pi_{\rm HMT}\frac{F_{107}}{F_{109}}
 -\frac{\pi_{\rm HMT}}{\varphi^2}
 \right|
 =
 1.9378907974286976068\ldots\times10^{-45}.
\]
Ces erreurs mesurent l'approximation finie du retour ; ce ne sont pas des ajustements de
\(\pi\) ni de \(\varphi\).

\paragraph{Lecteur cylindrique du rapport.}
Soient \(I_n^\pi=[\underline\pi_n,\overline\pi_n]\) et
\(I_n^\varphi=[\underline\varphi_n,\overline\varphi_n]\) les cylindres
positifs et emboîtés produits par les deux lecteurs HMT\@. L'opération sur les
intervalles
\[
 \mathcal Q(I,J)
 :=
 \left[
 \frac{\inf I}{(\sup J)^2},
 \frac{\sup I}{(\inf J)^2}
 \right]
\]
est le plus petit intervalle qui contient \(x/y^2\) pour tout
\((x,y)\in I\times J\). Elle préserve l'emboîtement : si les cylindres d'entrée sont
raffinés, \(\mathcal Q(I_n^\pi,I_n^\varphi)\) est également raffiné. Comme leurs
diamètres tendent vers zéro et que la limite de \(I_n^\varphi\) est positive,
\[
 \bigcap_n\mathcal Q(I_n^\pi,I_n^\varphi)
 =
 \left\{\frac{\pi_{\rm HMT}}{\varphi_{\rm HMT}^2}\right\}.
\]
Avec les douze blocs en base mille déjà publiés, l'intervalle quotient impose
trente-cinq décimales :
\[
 \frac{\pi_{\rm HMT}}{\varphi_{\rm HMT}^2}
 =
 1.19998161486432666111577775331680692\ldots.
\]
Cette construction conserve la provenance de chaque chiffre : le numérateur
provient du cylindre de clôture et le dénominateur de l'autoéchelle.

\paragraph{Séparation algébrique entre lecture aréale, lecture volumétrique et leurs codomaines respectifs}
Le rapport complet ne peut provenir de la seule algèbre rationnelle de
\(F_{\rm av}\). En effet, ses fonctions rationnelles spectrales appartiennent à
\(\mathbb Q(\sqrt5)\) ; on y trouve \(\varphi^{-2}\), mais pas le nombre
transcendant \(\pi\). Par conséquent, toute dérivation qui prétendrait
obtenir \(\pi/\varphi^2\) uniquement à partir de la matrice \(2\times2\) serait
mal typée. La donnée transcendante doit provenir du lecteur coinductif de
clôture, ou d'un cocycle de frontière équivalent démontré
indépendamment. Ce contrôle négatif fixe la répartition des tâches entre
la correspondance finie et le lecteur de nombre.

\subsubsection{Involution d'échange du quotient}
\label{sec:espejo-pi-phi-electron-hmt}

La correspondance admet deux orientations exactes d'une même loi. Si
\[
 \mathscr R(x,y)=\frac{x}{y^2},
 \qquad
 \mathsf J(x,y)=(y,x),
\]
alors \(\mathsf J^2=I\) et
\[
 \mathscr R(\pi,\varphi)=\frac{\pi}{\varphi^2},
 \qquad
 (\mathscr R\circ\mathsf J)(\pi,\varphi)
 =\frac{\varphi}{\pi^2}.
\]
On vérifie en outre les identités
\[
 \mathscr R(\pi,\varphi)\mathscr R(\varphi,\pi)
 =\frac1{\pi\varphi},
 \qquad
 \frac{\mathscr R(\pi,\varphi)}
 {\mathscr R(\varphi,\pi)}
 =\left(\frac{\pi}{\varphi}\right)^3.
\]
En écrivant
\(\mathcal B_{\rm av}^{(\pi)}=\mathscr R(\pi,\varphi)\), on obtient
également la reparamétrisation exacte
\[
 \frac{\varphi}{\pi^2}
 =
 \frac1{\varphi^3
 \bigl(\mathcal B_{\rm av}^{(\pi)}\bigr)^2},
 \qquad
 e^{\varphi/\pi^2}
 =
 \exp\!\left[
 \frac1{\varphi^3
 \bigl(\mathcal B_{\rm av}^{(\pi)}\bigr)^2}
 \right].
\]
Ainsi, le changement d'orientation entre clôture et croissance transporte le
rapport de frontière vers sa lecture échangée :
\[
 \frac{\pi}{\varphi^2}
 \xleftrightarrow{\ \mathsf J\ }
 \frac{\varphi}{\pi^2}.
\]
La première orientation lit la clôture aréale marquée sur la mémoire
volumétrique quadratique ; la seconde lit la croissance intérieure sur la
clôture quadratique. Cette correspondance est un théorème du lecteur HMT. La
mécanique dimensionnelle applique ensuite la seconde orientation dans la
construction électronique, sans faire de l'électron un antécédent de
\(\pi\), de \(\varphi\) ou de la correspondance.

\subsubsection{Principe d'Ambivalence des lecteurs aréal et volumétrique}
\label{sec:ambivalencia-grav-inercial-pi-phi2}

Le Principe d'Ambivalence affirme que la lecture aréale et la lecture
volumétrique sont deux fonctionnelles coordonnées, mais non interchangeables, d'un
même état avec mémoire. Leur coïncidence sur la diagonale est
\[
 \mathcal L_{\rm ar}^{\rm diag}(x)
 =\mathcal L_{\rm vol}^{\rm diag}(x).
\]
L'égalité diagonale n'efface pas l'ambivalence en dehors de celle-ci.

Sur la frontière aréale--volumétrique, les deux lecteurs se factorisent à travers
une même amplitude positive de mémoire \(\mathcal M(x)\) :
\[
 \mathcal L_{\rm vol}(x)=\mu_{\rm vol}\mathcal M(x),
 \qquad
 \mathcal L_{\rm ar}(x)=
 \pi_{\rm HMT}\mu_{\rm ar}\mathcal M(x).
\]
Alors le théorème de Parry et la marque de retour donnent
\[
 \boxed{\displaystyle
 \frac{\mathcal L_{\rm ar}(x)}
 {\mathcal L_{\rm vol}(x)}
 =\frac{\pi_{\rm HMT}}{\varphi^2}.}
\]
Avec la normalisation volumétrique, la différence relative des lecteurs est
\[
 \frac{\mathcal L_{\rm ar}-\mathcal L_{\rm vol}}
 {\mathcal L_{\rm vol}}
 =
 \frac{\pi_{\rm HMT}}{\varphi^2}-1
 =
 0.1999816148643266611\ldots.
\]
La lecture physique de cette correspondance est effectuée plus tard, dans le chapitre de
Cartan--Holst, après la construction de l'holonomie, de la cotétrade, de la connexion, de la torsion
et de la courbure. Ici, le résultat est entièrement HMT : un rapport exact entre
deux lecteurs de la même mémoire de feuillet.

Il convient de les distinguer de l'extracteur enrichi utilisé ensuite en
mécanique dimensionnelle. Là, une histoire complète conserve le comptage d'action
\(n_A\), les douze événements orientés \(e_0,\ldots,e_{11}\) et la cochaîne
torsionnelle au niveau du pas ; on en obtient
\[
 s_C=\sum_{i=0}^{5}(e_i+e_{i+6}),
 \qquad
 W_\pm=6n_A\pm s_C.
\]
Les canaux déjà engendrés par HMT satisfont alors l'identité exacte
\[
 \exp\!\left(n_AA_{\rm rad}+\frac{s_C}{6}C^\ast_{\rm rad}\right)
 =\left(q_+^{-W_+}q_-^{-W_-}\right)^{1/12}.
\]
Par conséquent, le passage chemin--caractère angulaire n'est pas absent : il s'effectue sur
le registre dodécaphasique enrichi. Ce que ce calcul de phase prouve à lui
seul, ce sont \(\Omega_{\rm sw},S,G\) ; il ne les remplace pas et ne permet pas de reconstruire le
registre complet après l'avoir projeté.

La mesure uniforme sur les \(104\,976\) présentations descend vers des atomes
de masses \(1/729\), \(1/243\) et \(1/54\), mais pas vers \(1/468\). Dans la
microfibre quintuple de \(\pi\), les deux probabilités sont
\[
 \tau_{468}(\mathscr R_\pi)=\frac5{468},
 \qquad
 F_\#\mu_{\rm seed}(\mathscr R_\pi)=\frac7{729}.
\]
Le relèvement moyenné comme le relèvement de phase réalisent la première : ils ne
transportent pas la mesure uniforme sur les germes, mais attribuent la même masse
totale à chaque classe observable. Le relèvement de phase conserve en outre
l'immobilité microscopique dans les pas neutres.

Pour relier le relèvement de phase aux histoires enrichies de
profondeur six, il faut un noyau \(\mathcal K_6\) sur
\(\mathcal H_6^{\rm cat}\times C_9\). Celui-ci descend vers
\(\widehat{\mathcal K}_{\rm ph}\) à travers
\(\rho_6\times\mathrm{id}_{C_9}\) si et seulement s'il satisfait, à chaque phase, la
condition d'agrégabilité forte
\[
 \sum_{\rho_6(g)=y}\mathcal K_6((h,r),(g,r+1))
 =L_{\tau(r)}(\rho_6(h),y),
\]
indépendamment de \(h\) dans chaque fibre de \(\rho_6\). Dans ce cas,
la composition avec \(F\) se projette sur \(\mathcal K_{\rm ph}\). La construction
d'une famille \((\mathcal K_m)_{m\geq6}\) qui satisfait cette identité et
commute avec toutes les troncatures est l'étape restante vers un
transfert de la limite inverse.

\section{Cocycle d'échelle et conservation cylindrique}
\label{sec:cociclo-parry-escala}

Soit \(r=(1,\varphi)\) un vecteur de Perron--Frobenius de
\(F_{\rm av}\). La matrice de transition de Parry est
\[
P_{ij}=\frac{(F_{\rm av})_{ij}r_j}{\varphi r_i}
=
\begin{pmatrix}
0&1\\
\varphi^{-2}&\varphi^{-1}
\end{pmatrix}.
\]
Pour une route admissible \(\gamma:i\to j\) de longueur \(n\),
\[
\mathbb P(\gamma\mid i)=\frac{r_j}{\varphi^nr_i},
\qquad
\boxed{\chi_P(\gamma)=\varphi^n\frac{r_i}{r_j}.}
\]
La télescopie démontre
\[
\chi_P(\gamma_2\circ\gamma_1)
=\chi_P(\gamma_2)\chi_P(\gamma_1).
\]
Sur une route fermée, \(\chi_P=\varphi^n\) ; son carré et son inverse
engendrent les combinaisons
\[
C_n=\frac{\chi_P^2+\chi_P^{-2}}2,\qquad
S_n=\frac{\chi_P^2-\chi_P^{-2}}2,
\]
qui satisfont
\[
x_{n+1}=3x_n-x_{n-1}.
\]

Comme \(F_{\rm av}\) est symétrique, soit
\[
\ell=\frac{r}{1+\varphi^2},\qquad \ell^{\mathsf T}r=1.
\]
La mesure d’un cylindre est
\[
\mu[x_0\cdots x_n]=\ell_{x_0}r_{x_n}\varphi^{-n}.
\]
Par conséquent,
\[
\boxed{
V_n(x)
=\frac{\mu[x_0]}{\mu[x_0\cdots x_n]}
=\varphi^n\frac{r_{x_0}}{r_{x_n}}
=\chi_P(x_0\cdots x_n).}
\]
Pour une échelle de dimension typée \(D>0\),
\[
a_n^{(D)}(x)=V_n(x)^{1/D}
\]
satisfait la loi de conservation cylindrique
\[
\boxed{
\bigl(a_n^{(D)}(x)\bigr)^D
\mu[x_0\cdots x_n]=\mu[x_0].}
\]

En rang spatial trois et sur des retours fermés, nous écrivons
\(a_n:=a_n^{(3)}\) et obtenons
\[
\boxed{
\frac{a_9}{a_0}=\varphi^3,\qquad
\frac{a_{27}}{a_0}=\varphi^9,\qquad
\frac{a_{108}}{a_0}=\varphi^{36}.}
\]
Si \(A_k:=a_{9k}\) pour \(k\geq1\), alors, pour \(k\geq2\),
\[
A_{k+1}-2A_k+A_{k-1}
=A_{k-1}(\varphi^3-1)^2>0.
\]
L'inégalité démontre la convexité en profondeur de retour ; une accélération
physique requiert en outre une carte de durée.

Le calcul premier de
\(\operatorname{Sym}^2(F_{\rm av})\) dans
\cref{sec:retorno-marcado-pi-phi2} démontre que son spectre est
\(\{\varphi^2,-1,\varphi^{-2}\}\). Nous conservons ici uniquement le corollaire
pour le cocycle d'échelle : sur une route fermée avec \(V_n=\varphi^n\), le
mode stable obéit à
\[
\boxed{M_n=\varphi^{-2n}=V_n^{-2}.}
\]
L'identification de ces trois espaces propres aux trois feuillets du TRIT
exige un opérateur d'entrelacement ; la coïncidence de cardinalité ne le remplace pas. Le
corollaire serait réfuté si la route fermée ne satisfaisait pas
\(V_n=\varphi^n\) ou si le mode stable n'avait pas pour valeur propre
\(\varphi^{-2}\).

\section{Théorème de prolongation dynamique : hypothèses globales et domaine certifié}

Les objets suivants du problème fini ont été construits et séparés
explicitement :
\[
 F,
 \qquad
 F_{\rm obs},
 \qquad
 \mathcal K_{\rm ph},
 \qquad
 \widehat{\mathcal K}_{\rm ph},
 \qquad
 K_{\rm cat},
 \qquad
 \tau_{468},
 \qquad
 F_{\rm av},
 \qquad
 \operatorname{Sym}^2(F_{\rm av}),
 \qquad
 D_\pi,
 \qquad
 \mathcal Q.
\]
Pour \(F_{\rm obs}\), le retour déterministe a été démontré et un
témoin explicite de non-descente produit a été exhibé. Pour \(\mathcal K_{\rm ph}\), on a
démontré la double stochasticité, l'irréductibilité, la période neuf, le
retour nonadique uniforme et l'unicité de
\(\pi_{\rm ph}\). Pour \(\widehat{\mathcal K}_{\rm ph}\), on a démontré
l'agrégabilité forte, la conservation microscopique de l'immobilité et la mesure
stationnaire compensée. \(K_{\rm cat}\) n'est plus une moyenne sans
provenance : c'est l'ombre d'un pas en moyennant la phase uniforme.
Le quotient de modes du TPK induit \(F_{\rm av}\) ; son élévation quadratique
sépare les modes \(\varphi^2,-1,\varphi^{-2}\). L'opérateur \(D_\pi\)
marque une seule fois le retour de frontière avec la clôture engendrée, et le
lecteur d'intervalles \(\mathcal Q\) conserve l'emboîtement des cylindres de
\(\pi\) et de \(\varphi\). Ces quatre objets construisent
\(\pi_{\rm HMT}/\varphi_{\rm HMT}^2\) sans confondre la fréquence temporelle,
la mesure uniforme et la mesure de Parry.

La force de cette induction est relative et concrète. Le mot TPK impose les
poids \(1/3\). L'interprétation APP de toutes les continuations compatibles,
avec le comptage uniforme, impose \(E_+\) et \(E_\times\). La trajectoire
déterministe \(F_{\rm obs}\) n'impose pas ces opérateurs et, de fait, ne les
produit pas. Ce n'est pas un défaut du résultat, mais la différence mathématique
entre l'évolution d'un chemin et le transfert d'un espace de chemins.

\ifdefined\HMTInternalTraceability
\begin{criterion}[Élévation enrichie du renversement]
\label{pf84:SYN31-RHO-LIFT}
On cherche un sous-domaine \(D_\rho\subset X_{\rm TPK}^{\rm enr}\) et une application
\[
 \Theta_\rho:D_\rho\longrightarrow X_{\rm TPK}^{\rm enr}
\]
qui relève le renversement entier \(27\mapsto72\) de
\cref{int:pf84:SYN30-RHO-DESC,int:pf84:SYN30-CROSS-WITNESS}. Elle doit commuter avec les
lecteurs résiduel, entier et généalogique, quantifier le quotient et la
retenue, et préserver l'orientation, la route, le retour et, lorsque le transport est
inversible, la monodromie, la région et la survie.
Ce n'est que si \(D_\rho\) est totale et invariante sous \(\Theta_\rho\) que l'on peut exiger
\(\Theta_\rho^2=I\) et la nommer involution.

Cette construction ni une démonstration d'involutivité n'existent encore. Le
programme est réfuté si l'application n'existe pas, si l'un des carrés de
lecture échoue, si l'on perd de la mémoire sans la quantifier ou si
\(\Theta_\rho^2\ne I\) dans un domaine où l'on revendique une involution.
\end{criterion}

\begin{criterion}[Élévation enrichie du carré]
\label{pf84:SYN31-SQ-LIFT}
On cherche un sous-domaine \(D_s\subset X_{\rm TPK}^{\rm enr}\) et une application
\[
 \Theta_s:D_s\longrightarrow X_{\rm TPK}^{\rm enr}
\]
qui relève \(s(n)=n^2\), en particulier \(27\mapsto729\), et commute avec les
lecteurs résiduel, entier et généalogique de
\cref{int:pf84:SYN30-SQ-DESC,int:pf84:SYN30-CROSS-WITNESS}. Elle doit quantifier le
quotient et la retenue et préserver l'orientation, la route, le retour et, lorsque le
transport est inversible, la monodromie, la région et la survie. Cette application n'est
pas involutive et on ne lui impose pas
\(\Theta_s^2=I\).

La démonstration reste absente. L'inexistence de \(\Theta_s\), l'échec
d'un carré de lecture ou la perte non contrôlée de l'une quelconque des
coordonnées enrichies réfuteraient le programme. L'ablation
\(81\mapsto72\) lors du changement de \(q^\times(9,9)\) est un contrôle indépendant :
ce n'est une branche ni de \(\Theta_\rho\) ni de \(\Theta_s\), et elle ne démontre aucune
des deux applications.
\end{criterion}
\fi

La conclusion finie est complète dans son domaine : le retour uniforme, la
mesure stationnaire, l'agrégabilité à profondeur six et la séparation
entre l'évolution d'une route et le transfert de l'espace des routes sont
démontrés par les objets construits. En particulier,
\(\widehat{\mathcal K}_{\rm ph}\) ne s'identifie pas à la chronologie de
\(F_{\rm obs}\), comme le montre \cref{prop:retorno-fobs-identidad}.
